\section{A reusable power-gap certificate for integral points}
\label{dio:section}

The Siegel packet contains an elementary square-completion certificate:
an identity $D^2f=A^2+R$ with $\deg R<\deg A$ bounds all integral points
of $y^2=f(x)$ \cite{repoSiegel}.  The following development extends that argument to every
integer exponent $m\geq2$, retains an equally simple height formula, and
gives an algorithm constructing the certificate on an explicitly described
family.  This resolves a class of integral-point problems; it does not
provide a general height bound for arbitrary curves.

All polynomials in this section are finite lists of integer coefficients
in ascending order, with trailing zeroes removed.  The zero polynomial is
the empty list.  Addition, multiplication, evaluation, degree, powers,
coefficient comparison, and the sum of absolute coefficients are fixed
primitive recursive functions on these lists.  No algorithm is asked to
evaluate an arbitrary primitive recursive program supplied as input.

\subsection{The discrete gap that replaces an analytic argument}

\begin{lemma}[Integer power gap]
\label{dio:gap}
For integers $a,b$ and $m\geq2$, if $b^m\ne a^m$, then
\[
             |b^m-a^m|\geq |a|^{m-1}.
\]
\end{lemma}
\begin{proof}
First take natural numbers $0\leq u<v$.  Monotonicity of powers gives
$u^{m-1}\leq v^{m-1}$, hence
\[
 u^m=u\,u^{m-1}\leq u\,v^{m-1}
                 \leq(v-1)v^{m-1}.
\]
Consequently $v^m-u^m\geq v^{m-1}$.

If $|a|\ne|b|$, apply this observation to their two absolute values.
The difference between their $m$th powers has absolute value at least
$\max(|a|,|b|)^{m-1}$, and the reverse triangle inequality gives
\[
 \bigl||b|^m-|a|^m\bigr|\leq |b^m-a^m|.
\]
This proves the claim in this case.  If $|a|=|b|$ and the two powers
are different, the powers are opposites, rather than equal.  They
are nonzero, and therefore $|a|\geq1$ and
\[
             |b^m-a^m|=2|a|^m\geq |a|^{m-1}.
\]
These alternatives exhaust the decidable comparisons of integers.
In particular the argument covers negative $a,b$ and both parities of $m$.
\end{proof}

The essential datum is discreteness: a nonzero integer has absolute value
at least one.  The lemma would be false with arbitrary rational $b$.
For instance a rational $b$ can approach a nonzero integer $a$ arbitrarily
closely without having $b^m=a^m$.  Thus denominator information must be
recorded explicitly when this lemma is used after a rational change of
variables.

\begin{lemma}[A lower bound for a polynomial]
\label{dio:polynomialbounds}
Let $A(X)=\sum_{i=0}^r a_iX^i$ be a nonconstant integer polynomial,
$a_r\ne0$, and put $L=\sum_{i<r}|a_i|$.  For an integer $x$ with
$X=|x|\geq1$,
\[
                |A(x)|\geq(X-L)X^{r-1}.
\]
If $G$ is a nonzero integer polynomial and $G(x)=0$, then
\[
             |x|\leq\sum_{i<\deg G}|g_i|
\]
when $G$ has positive degree; a nonzero constant has no root.
\end{lemma}
\begin{proof}
The leading coefficient has $|a_r|\geq1$.  The triangle inequality and
$X^i\leq X^{r-1}$ for $i<r$ give
\[
 |A(x)|\geq |a_r|X^r-\sum_{i<r}|a_i|X^i
           \geq X^r-LX^{r-1}.
\]
For the root assertion, separate $x=0$.  Otherwise the equation for a
root yields $X^s\leq(\sum_{i<s}|g_i|)X^{s-1}$, where $s=\deg G$.
Cancel the positive integer $X^{s-1}$.  No factorization or algebraic
closure is used.
\end{proof}

\begin{lemma}[Finite loss estimate]
\label{dio:loss}
If $X\geq L\geq0$ and $n\geq1$, then
\[
 (X-L)^n\geq X^n-nLX^{n-1}
            =(X-nL)X^{n-1}.
\]
\end{lemma}
\begin{proof}
Use the finite difference-of-powers identity
\[
 X^n-(X-L)^n
   =L\sum_{j=0}^{n-1}X^{n-1-j}(X-L)^j.
\]
Every summand is at most $X^{n-1}$, so the sum is at most
$nX^{n-1}$.  The identity itself follows by multiplying out and
cancelling adjacent terms.  A finite induction proves it uniformly in
$n$.
\end{proof}

\subsection{The certificate theorem}

\begin{definition}[Power-gap certificate]
\label{dio:certificate}
A certificate for the integral equation $y^m=f(x)$ is a tuple
$(m,D,f,A,R)$ of integers and coefficient lists satisfying
\[
 \begin{gathered}
 m\geq2,\qquad D\geq1,\qquad r=\deg A\geq1,\qquad R\ne0,\\
 D^m f=A^m+R,\qquad \deg R<(m-1)r.
 \end{gathered}
\]
The polynomial identity is a coefficientwise identity, not a finite
sample of values.  Each certificate condition is decidable by a bounded
calculation on the supplied finite lists.
\end{definition}

\begin{theorem}[Uniform integral-point bound]
\label{dio:height}
Given a power-gap certificate, define
\[
 \begin{split}
 L&=\sum_{i<r}|a_i|,\qquad C=\sum_i|r_i|,\\
 H&=(m-1)L+C,\qquad K=\sum_i|f_i|H^i.
 \end{split}
\]
Every integer solution of $y^m=f(x)$ satisfies
\[
                         |x|\leq H,\qquad |y|\leq K.
\]
The complete finite list of solutions is therefore primitive recursively
computable from the certificate.
\end{theorem}
\begin{proof}
Since $R\ne0$, $C\geq1$ and $H\geq1$.  Suppose $y^m=f(x)$ and put
$X=|x|$.  We prove that the decidable alternative $X>H$ is impossible.
First, $R(x)\ne0$.  A nonzero constant cannot vanish; if $R$ has positive
degree, Lemma~\ref{dio:polynomialbounds} bounds the absolute value of any
integer root by the sum of its lower coefficients, which is at most
$C\leq H$.

Put $n=m-1\geq1$.  The inequality $X>nL+C$ implies $X>L$.
Lemmas~\ref{dio:polynomialbounds} and~\ref{dio:loss} give
\begin{align*}
 |A(x)|^n
 &\geq (X-L)^n X^{n(r-1)}\\
 &\geq (X-nL)X^{n-1+n(r-1)}\\
 &> C X^{nr-1}.
\end{align*}
All quantities multiplied here are nonnegative, and $X^{nr-1}\geq1$.
On the other hand, every exponent in $R$ is at most $nr-1$, so
\[
                     |R(x)|\leq C X^{nr-1}<|A(x)|^{m-1}.
\]
Evaluating the certificate identity at $x$ and using $y^m=f(x)$ gives
\[
                    (Dy)^m-A(x)^m=R(x)\ne0.
\]
Lemma~\ref{dio:gap} gives the opposite inequality
$|R(x)|\geq|A(x)|^{m-1}$.  This contradiction rules out $X>H$ and
establishes $|x|\leq H$.

For the second coordinate, take absolute values in the equation:
\[
 |y|^m=|f(x)|\leq\sum_i|f_i||x|^i\leq K.
\]
Every natural number $t$ satisfies $t\leq t^m$ for $m\geq1$: the
case $t=0$ is immediate, and for $t\geq1$ use $t^{m-1}\geq1$.
Thus $|y|\leq K$.

Finally enumerate the integer rectangle
$[-H,H]\times[-K,K]$, retaining a pair exactly when its equation holds.
The equality test makes the list sound, and the proved bounds make it
complete.  These are bounded loops whose limits were computed before
the search.  The faster reference implementation scans only $x$ and
uses a bounded integer-root scan for $y$; its completeness is the same.
\end{proof}

For $m=2$, the formula is exactly $H=L+C$, recovering the square
certificate in the Siegel packet.  The extension to arbitrary $m$ does
not require a new analytic estimate or a real-root algorithm.  Its only
extra ingredient is the finite loss estimate.

\subsection{Constructing the certificates without an unbounded search}

\begin{theorem}[Polynomial power completion]
\label{dio:completion}
Let $m\geq2$, $r\geq1$, and let $f\in\mathbb Z[X]$ have degree $mr$
and leading coefficient $c^m$, where a nonzero integer $c$ is supplied.
A primitive recursive algorithm returns one of the following:
\begin{enumerate}
 \item a rational polynomial $Q$ with $f=Q^m$;
 \item a power-gap certificate for $y^m=f(x)$ and the complete finite
       solution list furnished by Theorem~\ref{dio:height}.
\end{enumerate}
The first alternative occurs exactly when $f$ is an $m$th power in
$\mathbb Q[X]$.
\end{theorem}
\begin{proof}
Start with $Q_0=cX^r$.  At stage $j$, for $1\leq j\leq r$, the
coefficients in degrees $r-1,\ldots,r-j+1$ have been assigned and all
lower coefficients are zero.  Let
\[
 d_j=[X^{mr-j}]Q_{j-1}^{m},\qquad
 b_{r-j}=\frac{f_{mr-j}-d_j}{m c^{m-1}},\qquad
 Q_j=Q_{j-1}+b_{r-j}X^{r-j}.
\]
The denominator is a specified nonzero integer.  In degree $mr-j$, the
new coefficient contributes precisely $m c^{m-1}b_{r-j}$: choose its
term in one of the $m$ factors and the leading term in every other
factor.  A term using it twice has degree at most $mr-2j<mr-j$.
Lower as-yet unassigned coefficients cannot contribute in this or a
higher degree.  Thus the new choice matches coefficient $mr-j$ and
preserves all earlier matches.

Finite induction on $j$ proves that $Q=Q_r$ has leading coefficient
$c$ and
\[
             [X^i](f-Q^m)=0\qquad((m-1)r\leq i\leq mr).
\]
Choose a positive common denominator $D$ for the coefficients of $Q$;
their denominator product suffices.  Then
\[
                    A=DQ,\qquad R=D^m f-A^m
\]
are integer polynomials.  If $R\ne0$, its degree is less than
$(m-1)r$, and all certificate conditions hold.  Its bounds and its
complete solution list follow from Theorem~\ref{dio:height}.

If $R=0$, report $f=Q^m$.  Conversely suppose $f=P^m$ for a rational
polynomial $P$.  Its degree is $r$.  Its leading coefficient equals $c$
when $m$ is odd, and equals $c$ or $-c$ when $m$ is even.  In the latter
case replace $P$ by $-P$ if necessary, preserving its $m$th power.
The coefficients of a root with leading coefficient $c$ must satisfy
the same triangular recurrence, one coefficient at a time.  Therefore
$P=Q$, and the algorithm reports the first alternative.

There are exactly $r$ rational coefficient steps.  Polynomial
convolution, integer powers, rational arithmetic, denominator products,
and the final zero test are all bounded finite computations.  The last
enumeration is bounded by the previously computed $H,K$.  In particular,
the construction is not a search for a height bound whose eventual
success has been postulated.
\end{proof}

For example, for $f=X^6+X^5+X+1$ and $m=3$ the construction begins
$Q=X^2+X/3-1/9$.  The reference code clears its denominators and checks
the resulting identity coefficient by coefficient.  The finite list
can be much more expensive to produce than its certificate; no claim
of polynomial-time complexity is intended.

\subsection{A complete new family}

\begin{theorem}[A uniform superelliptic classification]
\label{dio:family}
Let $m\geq3$ and $r\geq1$.  The integer points of
\[
                         y^m=x^{mr}+x+1
\]
are exactly
\[
 \begin{cases}
 \{(-1,(-1)^r),(0,1)\},&m\text{ odd},\\
 \{(-1,-1),(-1,1),(0,-1),(0,1)\},&m\text{ even}.
 \end{cases}
\]
\end{theorem}
\begin{proof}
Take $D=1$, $A=X^r$, and $R=X+1$.  The degree condition is
$1<(m-1)r$, which follows from $m\geq3,r\geq1$.  Here $L=0,C=2$,
so Theorem~\ref{dio:height} gives $|x|\leq2$.

For $x=-2$, the difference from $A(-2)^m$ is $-1$.  This is nonzero
and has absolute value smaller than
$|A(-2)|^{m-1}=2^{r(m-1)}\geq4$, contrary to Lemma~\ref{dio:gap}.
For $x=2$, the difference is $3$, again smaller than the same lower
bound $4$.  For $x=1$ the equation would be $y^m=3$.
If $|y|\leq1$, its power has absolute value at most one; if $|y|\geq2$,
its power has absolute value at least $2^3=8$.  Neither possibility
gives three.

The remaining values are $x=-1$ and $x=0$.  At $x=0$ the equation is
$y^m=1$.  It forces $|y|=1$ and permits only $y=1$ for odd $m$, or
$y=\pm1$ for even $m$.  At $x=-1$ it is
$y^m=(-1)^{mr}$, again forcing $|y|=1$.  For odd $m$ the sign is
$(-1)^r$, and for even $m$ both signs are allowed.  Direct substitution
verifies every point listed.
\end{proof}

This family includes cubic, quartic, and higher power equations of
arbitrarily large degree in $x$.  Its proof is uniform in both integer
parameters.  It extends the square family already treated in the
Siegel packet, rather than relabelling that earlier result.

\subsection{Transport with an explicit denominator bound}

The integer theorem can also be used for rational points when a bound
on denominators is part of the input.  This is a useful precise version
of an otherwise unsafe coordinate change.

\begin{corollary}[Bounded denominator transport]
\label{dio:denominators}
Let $(m,D,f,A,R)$ be a power-gap certificate, with $r,L,C$ as above,
and let $B\geq1$.  Put
\[
 \begin{split}
 H_B&=(m-1)B^{r+1}L+B^{m(r+1)}C,\\
 K_B&=B^m\sum_i |f_i|H_B^i.
 \end{split}
\]
Suppose integers $u,w$ and positive integers $v,z\leq B$ satisfy
\[
                         (w/z)^m=f(u/v).
\]
Then $|u|\leq H_B$ and $|w|\leq K_B$.  Reduced fractions are not
required.
\end{corollary}
\begin{proof}
The certificate identity implies $\deg f=mr$: the leading term of
$A^m$ cannot be cancelled by the lower-degree polynomial $R$.
Define integer polynomials in $U$ by
\[
 \begin{split}
 F_{v,z}(U)&=z^m v^{mr}f(U/v),\\
 A_{v,z}(U)&=zv^r A(U/v),\\
 R_{v,z}(U)&=z^m v^{mr}R(U/v).
 \end{split}
\]
Every exponent of $v$ after expansion is nonnegative.  The degrees of
$A_{v,z}$ and $R_{v,z}$ equal those of $A,R$, and $R_{v,z}\ne0$.
Furthermore
\[
       D^mF_{v,z}=A_{v,z}^m+R_{v,z},\qquad
       (v^r w)^m=F_{v,z}(u).
\]
This is another power-gap certificate.  Its lower-coefficient sum is
at most $B^{r+1}L$, and its remainder norm is at most
$B^{m(r+1)}C$.  Theorem~\ref{dio:height} therefore gives
$|u|\leq H_B$.  Finally
\[
 |w|^m=z^m|f(u/v)|
       \leq B^m\sum_i|f_i|H_B^i=K_B,
\]
and $|w|\leq|w|^m$ gives the second bound.
\end{proof}

The denominator parameter is substantive.  Letting $B$ depend on the
point gives no fixed finite set of rational points.  Thus this corollary
does not turn the integral theorem into a proof of Faltings' theorem.
It gives a complete, primitive recursive numerator bound once a single
denominator bound valid for the intended class of points has separately
been proved.  A birational transport must also supply the map, its
exceptional fibers, and its denominator estimate.

\subsection{Induction audit and Lean completion plan}

The certificate theorem has a particularly direct arithmetic translation:
\[
 \forall c,x,y\;\bigl(\operatorname{Valid}(c)=1\wedge
 y^{m(c)}=\operatorname{eval}(f(c),x)\bigr)
 \longrightarrow |x|\leq H(c)\wedge |y|\leq K(c).
\]
Here signed integer variables are natural-number codes.  Validity is a
primitive recursive Boolean function; the displayed formula uses only
first-order number quantifiers.  Its proof does not require a separate
sort of all polynomials or a real-number completion.

The finite sums, convolution identities, power monotonicity, and
difference-of-powers identity are proved by induction on their finite
loops.  Each invariant is an equality or inequality between primitive
recursive terms.  In the power-completion algorithm, the simultaneous
invariant says that all coefficients already processed agree with those
of $f$.  The coefficient index is bounded by $mr$; its bounded universal
quantifier is replaced by a Boolean fold.  Signed rational numerators
and positive denominators are finite codes, and the known divisor
$mc^{m-1}\ne0$ is carried through every step.  The final enumeration
uses the analogous Boolean prefix invariant: exactly the tested pairs
have been retained.  These are quantifier-free induction invariants
after the fixed finite data functions have been introduced.  Decidable
integer case distinctions do not use excluded middle for arbitrary
formulas.

The delivered \texttt{PowerGap.lean} is a compiled native Lean module
using only \texttt{Std}.  It proves natural-number upper and lower gap
lemmas, exclusions of small positive and negative perturbations, and
the full integer lemma, stated with positive exponent $n+1$:
\[
 b^{n+1}\ne a^{n+1}\;\Longrightarrow\;
 |a|^n\leq|b^{n+1}-a^{n+1}|.
\]
The two direct natural-number gap lemmas have no axiom dependencies.
The integer lemma and the two perturbation lemmas have only
\texttt{propext} and \texttt{Quot.sound} in their recorded native
dependency audit.  No new axiom, \texttt{Classical.choice}, or admitted
proof is used.  These remain native theorems, not embedded iPRA proof
objects.

The separate compiled file \texttt{PowerGapTargets.lean} supplies
executable coefficient-list definitions: normalization, coefficient
lookup, Horner evaluation, convolution, polynomial powers, coefficient
norms, a certificate record, the Boolean validity checker, and the two
bounds.  Its \texttt{Target} definitions name propositions and do not
assert them as theorems.  A complete implementation proceeds in this
order:
\begin{enumerate}
 \item Prove that normalization preserves coefficients and evaluation;
       prove evaluation respects convolution and powers.
 \item Prove the root bound and leading-term bound by finite coefficient
       induction, and the loss estimate by the finite power identity.
 \item Deduce the \texttt{OutsideBoxTarget}: beyond $H$, the remainder is
       nonzero and smaller than the allowed integer power gap.
 \item Combine the already proved integer gap with evaluated certificate
       identities to prove \texttt{HeightBoundTarget}; then prove exact
       bounded enumeration.
 \item Implement the rational coefficient recurrence, prove its bounded
       coefficient invariant, and prove the two output alternatives.
 \item Formalize the five possible values $-2,-1,0,1,2$ of $x$ to obtain
       \path{FamilyClassificationTarget}, with the explicit parity
       refinement in Theorem~\ref{dio:family}.
 \item Reify these fixed algorithms in the repository's primitive
       recursive syntax, replace native induction with the listed
       quantifier-free invariants, and construct the corresponding
       \texttt{Prf PRAAx} derivations.  Only after this step is the
       repository's existing theory-translation theorem applicable.
\end{enumerate}

This order exposes every remaining mathematical formalization obligation.
In particular the Boolean checker is not assumed sound merely because
it accepts examples, and Lean's proof of the integer gap is not silently
treated as an object-iPRA proof.

\subsection{What this resolves and what it leaves open}

The concrete advance is a certificate theorem and constructor for the
larger class $y^m=f(x)$ with $\deg f$ a positive multiple of $m$ and
leading coefficient an integer $m$th power, excluding exact polynomial
$m$th powers.  It includes a uniform family with a complete point
classification and a denominator-aware transport theorem.  Its data and
proof invariants fit the same finite arithmetic framework as the
existing square certificate.

Several other packets identify different missing information, and this
argument does not manufacture it.  The Mordell--Weil packet already
proves a finite height descent conditional on an actual rational list
covering the quotient; checking a list or a finite Selmer calculation
does not establish that covering assertion.  The Faltings packet
distinguishes a complete height bound from a bound on the number of
points.  The Roth packet similarly distinguishes a denominator cutoff
from a cardinality estimate or an empty finite window.  The
prime-distribution packet leaves the field-theoretic and analytic
contracts for Chebotarev, and the uniform parameter estimates needed
for Siegel--Walfisz and Bombieri--Vinogradov, unproved.  Recording these
requirements as input contracts is useful for composing proofs, but
does not discharge them \cite{repoMordell,repoFaltings,repoRoth,repoPrimes}.

For this reason the present generalization makes no claim to solve the
general effective Siegel, Faltings, Roth, or Mordell--Weil problems, and
does not assert that their unresolved iPRA targets are unprovable.  Its
contribution is an explicit additional family in which the missing
positive height information is actually constructed.
